\documentclass[11pt,a4paper]{article}

\newcommand{\doctitle}{Stress Testing}
\newcommand{\docsubtitle}{User Guide --- CDS-in-Stress, Flood Classifiers \& Portfolio Scenarios}
\newcommand{\docversion}{1.1}
\newcommand{\docdate}{29-October-2026}
\newcommand{\docauthor}{MKM Research Labs}

% MKM Research Labs — Shared Document Header
% Usage: % MKM Research Labs — Shared Document Header
% Usage: % MKM Research Labs — Shared Document Header
% Usage: \input{../shared/mkm_header}  (before \begin{document})
%
% Required commands before \input:
%   \newcommand{\doctitle}{...}
%   \newcommand{\docsubtitle}{...}
%   \newcommand{\docversion}{1.0}
%   \newcommand{\docdate}{13-February-2026}
%   \newcommand{\docauthor}{David K Kelly}

% ─── Packages ───────────────────────────────────────────────────────────────
\usepackage[utf8]{inputenc}
\usepackage{amsmath,amsfonts,amssymb}
\usepackage{listings}
\usepackage{color,xcolor}
\usepackage{graphicx}
\usepackage{hyperref}
\usepackage[toc,page]{appendix}
\usepackage{tabto}
\usepackage{float}
\usepackage{comment}
\usepackage[en-GB]{datetime2}
\usepackage{setspace}
\usepackage{subcaption}
\usepackage{booktabs}
\usepackage{longtable}
\usepackage{adjustbox}
\usepackage{enumitem}
\usepackage{fancyhdr}
\usepackage{lastpage}

% ─── Page geometry ──────────────────────────────────────────────────────────
\onehalfspacing
\usepackage{geometry}
\geometry{a4paper, total={170mm,247mm}, left=20mm, top=25mm, bottom=25mm}

% ─── Headers and footers ───────────────────────────────────────────────────
\pagestyle{fancy}
\fancyhf{}
\fancyhead[L]{\small\textcolor{gray}{MKM Research Labs}}
\fancyhead[R]{\small\textcolor{gray}{\doctitle{} --- v\docversion}}
\fancyfoot[C]{\small\textcolor{gray}{CONFIDENTIAL}}
\fancyfoot[R]{\small\textcolor{gray}{Page \thepage\ of \pageref{LastPage}}}
\renewcommand{\headrulewidth}{0.4pt}
\renewcommand{\footrulewidth}{0.4pt}

% ─── Colours ────────────────────────────────────────────────────────────────
\definecolor{mkmblue}{HTML}{1565C0}
\definecolor{codegreen}{rgb}{0,0.6,0}
\definecolor{codegray}{rgb}{0.5,0.5,0.5}
\definecolor{codepurple}{rgb}{0.58,0,0.82}
\definecolor{backcolour}{rgb}{0.95,0.95,0.92}

% ─── Hyperref setup ─────────────────────────────────────────────────────────
\hypersetup{
    colorlinks=true,
    linkcolor=mkmblue,
    citecolor=mkmblue,
    urlcolor=mkmblue
}

% ─── Code listings ──────────────────────────────────────────────────────────
\lstdefinestyle{mystyle}{
    backgroundcolor=\color{backcolour},
    commentstyle=\color{codegreen},
    keywordstyle=\color{magenta},
    numberstyle=\tiny\color{codegray},
    stringstyle=\color{codepurple},
    basicstyle=\ttfamily\footnotesize,
    breakatwhitespace=false,
    breaklines=true,
    captionpos=b,
    keepspaces=true,
    numbers=left,
    numbersep=5pt,
    showspaces=false,
    showstringspaces=false,
    showtabs=false,
    tabsize=2
}
\lstset{style=mystyle}

% ─── Section formatting ────────────────────────────────────────────────────
\setcounter{tocdepth}{3}
\setcounter{secnumdepth}{3}

% ─── Shared macros ──────────────────────────────────────────────────────────
\newcommand{\mkmtitlepage}{%
    \begin{titlepage}
    \begin{center}
    \vspace*{1cm}
    \begin{figure}[H]\centering
        \includegraphics[width=0.4\textwidth]{../shared/Logos/MKM_logo.png}
    \end{figure}
    \vspace{1.5cm}
    {\Huge\bfseries\textcolor{mkmblue}{\doctitle}\par}
    \vspace{0.5cm}
    {\Large\textcolor{gray}{\docsubtitle}\par}
    \vspace{2cm}
    {\large MKM Research Labs\par}
    \vspace{0.5cm}
    {\large Version \docversion{} --- \docdate\par}
    \vspace{0.5cm}
    {\large \docauthor\par}
    \vspace{2cm}
    \textcolor{red}{\textbf{CONFIDENTIAL --- PROPRIETARY}}
    \end{center}
    \end{titlepage}
}

\newcommand{\mkmlegalpage}{%
    \clearpage
    \thispagestyle{empty}
    \begin{center}
    \large\textbf{Legal Notice}
    \vspace{1em}

    \textbf{PROPRIETARY AND CONFIDENTIAL}
    \end{center}

    \noindent This document, including all algorithms, models, methodology, formulas,
    calculations, and intellectual concepts contained herein, constitutes the exclusive
    intellectual property and confidential trade secrets of MKM Research Labs (``MKM'').

    \vspace{0.5em}
    \noindent Any usage, reproduction, distribution, modification, or disclosure of this
    document or any portion thereof, without the express prior written authorisation from
    MKM Research Labs is strictly prohibited and constitutes an infringement of intellectual
    property rights.

    \vspace{0.5em}
    \noindent The document is provided on an ``as is'' basis. MKM Research Labs makes no
    representations or warranties of any kind, express or implied, regarding its accuracy,
    reliability, or suitability for any particular purpose.

    \vspace{1em}
    \begin{center}
    All rights reserved. \copyright\ 2021--2026 MKM Research Labs.

    \vspace{0.5em}
    Governed by the laws of the United Kingdom.
    \end{center}
    \clearpage
}

\newcommand{\mkmhistorypage}[1]{%
    \section{Document History}
    \begin{table}[ht]
        \centering
        \begin{tabular}{l|l|l|l}
            \textbf{Date} & \textbf{Version} & \textbf{Description} & \textbf{Author} \\
            \hline
            #1
        \end{tabular}
        \label{tab:revision_history}
    \end{table}
    \clearpage
}

% ─── Bank validation protocol section commands ──────────────────────────────
% Use these to ensure every model doc has the required sections.
% They are just styled \section calls — content is provided by each doc.

\newcommand{\mkmexecsummary}{\section{Executive Summary}}
\newcommand{\mkmpurpose}{\section{Model Purpose and Scope}}
\newcommand{\mkmframework}{\section{Mathematical Framework}}
\newcommand{\mkminputs}{\section{Input Parameters and Calibration}}
\newcommand{\mkmimplementation}{\section{Implementation Details}}
\newcommand{\mkmsensitivity}{\section{Sensitivity Analysis}}
\newcommand{\mkmvalidation}{\section{Validation and Backtesting}}
\newcommand{\mkmtestresults}{\subsection{Automated Test Results}\input{test_results}}
\newcommand{\mkmlimitations}{\section{Model Limitations and Known Weaknesses}}

\newcommand{\mkmgovernance}{\section{Governance Validation Questions}
\noindent The following nine questions are mandated by Chapter~5 of the
\emph{Model Risk Governance for Vendors} handbook and must be explicitly
addressed for every model in the inventory. Response status is tracked in the
Model Governance Dashboard (MRC portal). \par\medskip}

\newcommand{\mkmvqitem}[3]{\subsection{Q#1 --- #2}
\noindent\textit{``#3''}\par\medskip\noindent}

\newcommand{\mkmchangelog}{\section{Change History}}
  (before \begin{document})
%
% Required commands before \input:
%   \newcommand{\doctitle}{...}
%   \newcommand{\docsubtitle}{...}
%   \newcommand{\docversion}{1.0}
%   \newcommand{\docdate}{13-February-2026}
%   \newcommand{\docauthor}{David K Kelly}

% ─── Packages ───────────────────────────────────────────────────────────────
\usepackage[utf8]{inputenc}
\usepackage{amsmath,amsfonts,amssymb}
\usepackage{listings}
\usepackage{color,xcolor}
\usepackage{graphicx}
\usepackage{hyperref}
\usepackage[toc,page]{appendix}
\usepackage{tabto}
\usepackage{float}
\usepackage{comment}
\usepackage[en-GB]{datetime2}
\usepackage{setspace}
\usepackage{subcaption}
\usepackage{booktabs}
\usepackage{longtable}
\usepackage{adjustbox}
\usepackage{enumitem}
\usepackage{fancyhdr}
\usepackage{lastpage}

% ─── Page geometry ──────────────────────────────────────────────────────────
\onehalfspacing
\usepackage{geometry}
\geometry{a4paper, total={170mm,247mm}, left=20mm, top=25mm, bottom=25mm}

% ─── Headers and footers ───────────────────────────────────────────────────
\pagestyle{fancy}
\fancyhf{}
\fancyhead[L]{\small\textcolor{gray}{MKM Research Labs}}
\fancyhead[R]{\small\textcolor{gray}{\doctitle{} --- v\docversion}}
\fancyfoot[C]{\small\textcolor{gray}{CONFIDENTIAL}}
\fancyfoot[R]{\small\textcolor{gray}{Page \thepage\ of \pageref{LastPage}}}
\renewcommand{\headrulewidth}{0.4pt}
\renewcommand{\footrulewidth}{0.4pt}

% ─── Colours ────────────────────────────────────────────────────────────────
\definecolor{mkmblue}{HTML}{1565C0}
\definecolor{codegreen}{rgb}{0,0.6,0}
\definecolor{codegray}{rgb}{0.5,0.5,0.5}
\definecolor{codepurple}{rgb}{0.58,0,0.82}
\definecolor{backcolour}{rgb}{0.95,0.95,0.92}

% ─── Hyperref setup ─────────────────────────────────────────────────────────
\hypersetup{
    colorlinks=true,
    linkcolor=mkmblue,
    citecolor=mkmblue,
    urlcolor=mkmblue
}

% ─── Code listings ──────────────────────────────────────────────────────────
\lstdefinestyle{mystyle}{
    backgroundcolor=\color{backcolour},
    commentstyle=\color{codegreen},
    keywordstyle=\color{magenta},
    numberstyle=\tiny\color{codegray},
    stringstyle=\color{codepurple},
    basicstyle=\ttfamily\footnotesize,
    breakatwhitespace=false,
    breaklines=true,
    captionpos=b,
    keepspaces=true,
    numbers=left,
    numbersep=5pt,
    showspaces=false,
    showstringspaces=false,
    showtabs=false,
    tabsize=2
}
\lstset{style=mystyle}

% ─── Section formatting ────────────────────────────────────────────────────
\setcounter{tocdepth}{3}
\setcounter{secnumdepth}{3}

% ─── Shared macros ──────────────────────────────────────────────────────────
\newcommand{\mkmtitlepage}{%
    \begin{titlepage}
    \begin{center}
    \vspace*{1cm}
    \begin{figure}[H]\centering
        \includegraphics[width=0.4\textwidth]{../shared/Logos/MKM_logo.png}
    \end{figure}
    \vspace{1.5cm}
    {\Huge\bfseries\textcolor{mkmblue}{\doctitle}\par}
    \vspace{0.5cm}
    {\Large\textcolor{gray}{\docsubtitle}\par}
    \vspace{2cm}
    {\large MKM Research Labs\par}
    \vspace{0.5cm}
    {\large Version \docversion{} --- \docdate\par}
    \vspace{0.5cm}
    {\large \docauthor\par}
    \vspace{2cm}
    \textcolor{red}{\textbf{CONFIDENTIAL --- PROPRIETARY}}
    \end{center}
    \end{titlepage}
}

\newcommand{\mkmlegalpage}{%
    \clearpage
    \thispagestyle{empty}
    \begin{center}
    \large\textbf{Legal Notice}
    \vspace{1em}

    \textbf{PROPRIETARY AND CONFIDENTIAL}
    \end{center}

    \noindent This document, including all algorithms, models, methodology, formulas,
    calculations, and intellectual concepts contained herein, constitutes the exclusive
    intellectual property and confidential trade secrets of MKM Research Labs (``MKM'').

    \vspace{0.5em}
    \noindent Any usage, reproduction, distribution, modification, or disclosure of this
    document or any portion thereof, without the express prior written authorisation from
    MKM Research Labs is strictly prohibited and constitutes an infringement of intellectual
    property rights.

    \vspace{0.5em}
    \noindent The document is provided on an ``as is'' basis. MKM Research Labs makes no
    representations or warranties of any kind, express or implied, regarding its accuracy,
    reliability, or suitability for any particular purpose.

    \vspace{1em}
    \begin{center}
    All rights reserved. \copyright\ 2021--2026 MKM Research Labs.

    \vspace{0.5em}
    Governed by the laws of the United Kingdom.
    \end{center}
    \clearpage
}

\newcommand{\mkmhistorypage}[1]{%
    \section{Document History}
    \begin{table}[ht]
        \centering
        \begin{tabular}{l|l|l|l}
            \textbf{Date} & \textbf{Version} & \textbf{Description} & \textbf{Author} \\
            \hline
            #1
        \end{tabular}
        \label{tab:revision_history}
    \end{table}
    \clearpage
}

% ─── Bank validation protocol section commands ──────────────────────────────
% Use these to ensure every model doc has the required sections.
% They are just styled \section calls — content is provided by each doc.

\newcommand{\mkmexecsummary}{\section{Executive Summary}}
\newcommand{\mkmpurpose}{\section{Model Purpose and Scope}}
\newcommand{\mkmframework}{\section{Mathematical Framework}}
\newcommand{\mkminputs}{\section{Input Parameters and Calibration}}
\newcommand{\mkmimplementation}{\section{Implementation Details}}
\newcommand{\mkmsensitivity}{\section{Sensitivity Analysis}}
\newcommand{\mkmvalidation}{\section{Validation and Backtesting}}
\newcommand{\mkmtestresults}{\subsection{Automated Test Results}\input{test_results}}
\newcommand{\mkmlimitations}{\section{Model Limitations and Known Weaknesses}}

\newcommand{\mkmgovernance}{\section{Governance Validation Questions}
\noindent The following nine questions are mandated by Chapter~5 of the
\emph{Model Risk Governance for Vendors} handbook and must be explicitly
addressed for every model in the inventory. Response status is tracked in the
Model Governance Dashboard (MRC portal). \par\medskip}

\newcommand{\mkmvqitem}[3]{\subsection{Q#1 --- #2}
\noindent\textit{``#3''}\par\medskip\noindent}

\newcommand{\mkmchangelog}{\section{Change History}}
  (before \begin{document})
%
% Required commands before \input:
%   \newcommand{\doctitle}{...}
%   \newcommand{\docsubtitle}{...}
%   \newcommand{\docversion}{1.0}
%   \newcommand{\docdate}{13-February-2026}
%   \newcommand{\docauthor}{David K Kelly}

% ─── Packages ───────────────────────────────────────────────────────────────
\usepackage[utf8]{inputenc}
\usepackage{amsmath,amsfonts,amssymb}
\usepackage{listings}
\usepackage{color,xcolor}
\usepackage{graphicx}
\usepackage{hyperref}
\usepackage[toc,page]{appendix}
\usepackage{tabto}
\usepackage{float}
\usepackage{comment}
\usepackage[en-GB]{datetime2}
\usepackage{setspace}
\usepackage{subcaption}
\usepackage{booktabs}
\usepackage{longtable}
\usepackage{adjustbox}
\usepackage{enumitem}
\usepackage{fancyhdr}
\usepackage{lastpage}

% ─── Page geometry ──────────────────────────────────────────────────────────
\onehalfspacing
\usepackage{geometry}
\geometry{a4paper, total={170mm,247mm}, left=20mm, top=25mm, bottom=25mm}

% ─── Headers and footers ───────────────────────────────────────────────────
\pagestyle{fancy}
\fancyhf{}
\fancyhead[L]{\small\textcolor{gray}{MKM Research Labs}}
\fancyhead[R]{\small\textcolor{gray}{\doctitle{} --- v\docversion}}
\fancyfoot[C]{\small\textcolor{gray}{CONFIDENTIAL}}
\fancyfoot[R]{\small\textcolor{gray}{Page \thepage\ of \pageref{LastPage}}}
\renewcommand{\headrulewidth}{0.4pt}
\renewcommand{\footrulewidth}{0.4pt}

% ─── Colours ────────────────────────────────────────────────────────────────
\definecolor{mkmblue}{HTML}{1565C0}
\definecolor{codegreen}{rgb}{0,0.6,0}
\definecolor{codegray}{rgb}{0.5,0.5,0.5}
\definecolor{codepurple}{rgb}{0.58,0,0.82}
\definecolor{backcolour}{rgb}{0.95,0.95,0.92}

% ─── Hyperref setup ─────────────────────────────────────────────────────────
\hypersetup{
    colorlinks=true,
    linkcolor=mkmblue,
    citecolor=mkmblue,
    urlcolor=mkmblue
}

% ─── Code listings ──────────────────────────────────────────────────────────
\lstdefinestyle{mystyle}{
    backgroundcolor=\color{backcolour},
    commentstyle=\color{codegreen},
    keywordstyle=\color{magenta},
    numberstyle=\tiny\color{codegray},
    stringstyle=\color{codepurple},
    basicstyle=\ttfamily\footnotesize,
    breakatwhitespace=false,
    breaklines=true,
    captionpos=b,
    keepspaces=true,
    numbers=left,
    numbersep=5pt,
    showspaces=false,
    showstringspaces=false,
    showtabs=false,
    tabsize=2
}
\lstset{style=mystyle}

% ─── Section formatting ────────────────────────────────────────────────────
\setcounter{tocdepth}{3}
\setcounter{secnumdepth}{3}

% ─── Shared macros ──────────────────────────────────────────────────────────
\newcommand{\mkmtitlepage}{%
    \begin{titlepage}
    \begin{center}
    \vspace*{1cm}
    \begin{figure}[H]\centering
        \includegraphics[width=0.4\textwidth]{../shared/Logos/MKM_logo.png}
    \end{figure}
    \vspace{1.5cm}
    {\Huge\bfseries\textcolor{mkmblue}{\doctitle}\par}
    \vspace{0.5cm}
    {\Large\textcolor{gray}{\docsubtitle}\par}
    \vspace{2cm}
    {\large MKM Research Labs\par}
    \vspace{0.5cm}
    {\large Version \docversion{} --- \docdate\par}
    \vspace{0.5cm}
    {\large \docauthor\par}
    \vspace{2cm}
    \textcolor{red}{\textbf{CONFIDENTIAL --- PROPRIETARY}}
    \end{center}
    \end{titlepage}
}

\newcommand{\mkmlegalpage}{%
    \clearpage
    \thispagestyle{empty}
    \begin{center}
    \large\textbf{Legal Notice}
    \vspace{1em}

    \textbf{PROPRIETARY AND CONFIDENTIAL}
    \end{center}

    \noindent This document, including all algorithms, models, methodology, formulas,
    calculations, and intellectual concepts contained herein, constitutes the exclusive
    intellectual property and confidential trade secrets of MKM Research Labs (``MKM'').

    \vspace{0.5em}
    \noindent Any usage, reproduction, distribution, modification, or disclosure of this
    document or any portion thereof, without the express prior written authorisation from
    MKM Research Labs is strictly prohibited and constitutes an infringement of intellectual
    property rights.

    \vspace{0.5em}
    \noindent The document is provided on an ``as is'' basis. MKM Research Labs makes no
    representations or warranties of any kind, express or implied, regarding its accuracy,
    reliability, or suitability for any particular purpose.

    \vspace{1em}
    \begin{center}
    All rights reserved. \copyright\ 2021--2026 MKM Research Labs.

    \vspace{0.5em}
    Governed by the laws of the United Kingdom.
    \end{center}
    \clearpage
}

\newcommand{\mkmhistorypage}[1]{%
    \section{Document History}
    \begin{table}[ht]
        \centering
        \begin{tabular}{l|l|l|l}
            \textbf{Date} & \textbf{Version} & \textbf{Description} & \textbf{Author} \\
            \hline
            #1
        \end{tabular}
        \label{tab:revision_history}
    \end{table}
    \clearpage
}

% ─── Bank validation protocol section commands ──────────────────────────────
% Use these to ensure every model doc has the required sections.
% They are just styled \section calls — content is provided by each doc.

\newcommand{\mkmexecsummary}{\section{Executive Summary}}
\newcommand{\mkmpurpose}{\section{Model Purpose and Scope}}
\newcommand{\mkmframework}{\section{Mathematical Framework}}
\newcommand{\mkminputs}{\section{Input Parameters and Calibration}}
\newcommand{\mkmimplementation}{\section{Implementation Details}}
\newcommand{\mkmsensitivity}{\section{Sensitivity Analysis}}
\newcommand{\mkmvalidation}{\section{Validation and Backtesting}}
\newcommand{\mkmtestresults}{\subsection{Automated Test Results}\input{test_results}}
\newcommand{\mkmlimitations}{\section{Model Limitations and Known Weaknesses}}

\newcommand{\mkmgovernance}{\section{Governance Validation Questions}
\noindent The following nine questions are mandated by Chapter~5 of the
\emph{Model Risk Governance for Vendors} handbook and must be explicitly
addressed for every model in the inventory. Response status is tracked in the
Model Governance Dashboard (MRC portal). \par\medskip}

\newcommand{\mkmvqitem}[3]{\subsection{Q#1 --- #2}
\noindent\textit{``#3''}\par\medskip\noindent}

\newcommand{\mkmchangelog}{\section{Change History}}


\begin{document}

\mkmtitlepage
\mkmlegalpage

\tableofcontents
\clearpage


% =========================================================================
\section{Introduction}
% =========================================================================

Stress testing answers the question: \emph{what would happen to our
PRS positions if a specific storm hit the catchment?}

The platform uses its catalogue of 20,000 synthetic storm sequences to
run hourly scenarios.  For each hour of a 168-hour storm, the system
computes the gauge water level, the probability of flooding, and the
resulting P\&L impact on every affected trade.  This produces a
time-resolved view of how the portfolio behaves under stress.

\subsection{CDS-in-Stress Concept}

In normal mark-to-market, PRS positions are valued using CDS-style
analytics (hazard rate $\to$ fair spread $\to$ MTM).  Under stress,
the position transitions from a credit instrument to a
\textbf{cash-settled physical payout}:

\[
  \text{cash price} = N \times P(\text{flood}) \times d
\]

\noindent where $N$ is the notional, $P(\text{flood})$ is the real-time
flood probability, and $d$ is the direction (+1 payer, $-$1 receiver).
The stress P\&L is the difference between this cash price and the
normal-market MTM.


% =========================================================================
\section{Gauge-Level Stress}
% =========================================================================

\subsection{Workflow}

\begin{enumerate}
  \item \textbf{Select a gauge} from the dropdown.  Only gauges with
        open trades are shown.
  \item \textbf{Select a storm} from the catalogue.  The list shows
        storms that breached alert level at this gauge, sorted by
        severity (catastrophic first).
  \item \textbf{Run the scenario.}  The server computes the hourly
        P\&L trajectory over 168 hours.
\end{enumerate}

\subsection{What the Charts Show}

The stress tab displays three sub-charts:

\begin{description}[leftmargin=2em]
  \item[Flood Probability] A line chart of $P(\text{flood})$ vs.\
        hour.  Shows how the probability rises as water level
        approaches the trigger threshold, peaks, and then recedes.
  \item[Stress P\&L] A line chart of portfolio stress P\&L vs.\ hour.
        Shows the accumulated cash-price impact of the storm.
  \item[Probability Surface] A 3D surface of $P(\text{flood})$ vs.\
        water level vs.\ hour.  Useful for understanding how
        sensitivity changes at different water levels.
\end{description}


% =========================================================================
\section{The Hourly Scenario}
% =========================================================================

For each hour $h$ in the 168-hour window:

\begin{enumerate}
  \item \textbf{Water level} $w_h$ is read from the storm's gauge
        response hydrograph (pre-computed from the storm catalogue).
  \item \textbf{Flood probability} $p_h = f(w_h, h)$ is computed
        using either a trained GBM classifier or the analytical flood
        polynomial.
  \item \textbf{Cash price} per trade: $C_h = N \times p_h \times d$.
  \item \textbf{Stress P\&L} per trade: $\Pi_h = C_h - \text{MTM}$.
  \item \textbf{Portfolio stress P\&L}: $\sum \Pi_h$ across all
        affected trades.
\end{enumerate}

\subsection{Knock-Out Behaviour}

When the water level exceeds a trade's trigger threshold:
\begin{itemize}
  \item The trade ``knocks out'' --- the protection buyer receives
        full notional.
  \item $P(\text{flood}) = 1.0$ from the knock-out hour onward.
  \item Stress P\&L jumps to $N - \text{MTM}$ (full notional minus
        the original mark).
  \item The trade is no longer re-priced in subsequent hours.
\end{itemize}

\noindent In the chart, the P\&L line flattens after knock-out because
no further P\&L change can occur.

\subsection{Hydrograph Source}

The hourly water levels come from the gauge response stored in
\texttt{gaugets/\{gauge\_id\}.json}.  The storm's \texttt{peak\_level\_m}
and duration are used to scale the hydrograph template.  If the gauge
time series is not available, the system falls back to a flat hydrograph
(constant water level for 168 hours).


% =========================================================================
\section{Flood Probability Models}
% =========================================================================

Two models can compute $P(\text{flood})$.  The classifier is preferred
when available because it captures non-linear effects.

\subsection{GBM Classifier (Preferred)}

A \textbf{Gradient Boosting Machine} trained per gauge using historical
storm simulations.  The model takes four features:

\begin{center}
\begin{tabular}{lp{8cm}}
\toprule
\textbf{Feature} & \textbf{Description} \\
\midrule
\texttt{log\_h\_s} & Log of water level relative to severe threshold:
  $\ln(w / w_{\text{severe}} + \varepsilon)$ \\
\texttt{log\_t\_end} & Log of time remaining in the 168-hour window:
  $\ln(168 - h + \varepsilon)$ \\
\texttt{delta\_log\_h} & First derivative of log water level
  (rate of rise/fall) \\
\texttt{delta2\_log\_h} & Second derivative (acceleration of
  rise/fall) \\
\bottomrule
\end{tabular}
\end{center}

\noindent The classifier outputs a probability in $[0, 1]$.  It is
more accurate than the polynomial model because it learns the
non-linear relationship between water level dynamics and flood
outcomes from the Monte Carlo data.

\subsection{Analytical Flood Polynomial (Fallback)}

When no classifier is trained for a gauge, the system uses a
polynomial approximation:

\[
  P(\text{flood}) = \sigma\!\left(
    a_0 + a_1 \frac{w}{w_s} + a_2 \left(\frac{w}{w_s}\right)^2
    + a_3 \frac{h}{168}
  \right)
\]

\noindent where $\sigma$ is the logistic sigmoid and the coefficients
are calibrated to the Thames.  This model is smooth but cannot capture
the non-linear dynamics learned by the classifier.

\subsection{When to Use Which}

\begin{center}
\begin{tabular}{lll}
\toprule
& \textbf{GBM Classifier} & \textbf{Flood Polynomial} \\
\midrule
Accuracy     & Higher (trained on data) & Lower (generic fit) \\
Availability & Requires training        & Always available \\
Speed        & Fast (tree lookup)       & Instant (formula) \\
Per-gauge    & Yes (gauge-specific)     & No (one model for all) \\
\bottomrule
\end{tabular}
\end{center}


% =========================================================================
\section{Portfolio-Level Stress}
% =========================================================================

The Portfolio Stress tab runs a single storm across \textbf{all gauges}
simultaneously, computing the aggregate impact.

\subsection{Workflow}

\begin{enumerate}
  \item The tab lists all storms in the catalogue, sorted by the
        number of gauges that breached severe (most impactful first).
  \item Select a storm.
  \item The server runs the scenario for every gauge that has open
        trades and returns per-gauge results.
\end{enumerate}

\subsection{Output}

For each gauge, the response includes:
\begin{itemize}
  \item Peak water level and threshold classification (alert,
        warning, severe, or clean).
  \item $P(\text{flood})$ at peak.
  \item Stress P\&L contribution.
  \item Number of affected trades.
  \item Full 168-hour hydrograph.
\end{itemize}

\noindent The portfolio summary shows:
\begin{itemize}
  \item Total stress P\&L across all gauges.
  \item Number of gauges at each severity level.
  \item Worst-hit gauge and its contribution.
\end{itemize}

\subsection{Drill-Down}

Clicking a gauge in the portfolio stress results switches to the
Stress tab (single gauge) and pre-selects that gauge and storm,
allowing detailed hourly analysis.


% =========================================================================
\section{Probability Surface}
% =========================================================================

The probability surface is a 3D visualisation of $P(\text{flood})$ as
a function of water level and hour.  It is generated by evaluating the
classifier (or polynomial) across a grid of water levels and time
points.

\subsection{How to Read It}

\begin{itemize}
  \item The x-axis is time (0--168 hours).
  \item The y-axis is water level (metres).
  \item The z-axis (colour) is $P(\text{flood})$.
  \item The actual storm trajectory is a path across this surface.
\end{itemize}

\noindent The surface reveals the classifier's learned behaviour:
\begin{itemize}
  \item At low water levels, $P(\text{flood}) \approx 0$ regardless
        of hour.
  \item At high water levels, $P(\text{flood}) \approx 1$ regardless
        of hour.
  \item The transition zone (e.g.\ 0.2--0.8 probability) is where
        the storm's trajectory determines the outcome.
  \item Late in the window (hour 140+), the probability decreases
        because there is insufficient time for a full flood event to
        develop.
\end{itemize}


% =========================================================================
\section{Classifier Training}
% =========================================================================

\subsection{Training a Classifier}

To train a GBM classifier for a gauge:
\begin{enumerate}
  \item Navigate to the Classifiers tab.
  \item Select the gauge.
  \item Click \textbf{Train}.  The system builds a feature matrix from
        all storm scenarios at this gauge and trains a GBM model.
  \item Training typically takes 10--60 seconds per gauge.
  \item The trained model is saved as a \texttt{.joblib} file in
        \texttt{classifiers/\{gauge\_id\}.joblib}.
\end{enumerate}

\subsection{Batch Training}

The \textbf{Train All} button trains classifiers for all gauges with
sufficient historical data.  This is an asynchronous job; progress is
shown in the UI.

\subsection{AUC Metric}

After training, the classifier's quality is measured by the
\textbf{Area Under the ROC Curve} (AUC):
\begin{description}[leftmargin=2em]
  \item[$> 0.95$] Excellent --- very reliable flood probability
        estimates.
  \item[$0.85$--$0.95$] Good --- reliable for stress testing.
  \item[$< 0.85$] Fair --- consider retraining with more data or
        using the polynomial fallback.
\end{description}

\subsection{When to Retrain}

Classifiers should be retrained when:
\begin{itemize}
  \item The storm catalogue is regenerated (new storm control
        parameters).
  \item Gauge hazard data is updated.
  \item The AUC drops below 0.85 (indicates the model is stale).
\end{itemize}


% =========================================================================
\section{Model Selection: Challenger Comparison}
\label{sec:challenger}
% =========================================================================

\subsection{Open Question}

\begin{quote}
\emph{Should an alternative classifier, such as a Random Forest, be
considered in place of the GBM?  What evidence shows that the GBM is
the most accurate and most stable choice for $P(\text{flood})$?}
\end{quote}

\noindent The GBM was selected on the basis of in-sample AUC.  That is
not sufficient evidence for model selection: no challenger has been
fitted on the same data, and the AUC is measured on a split that
overstates out-of-sample performance (Section~\ref{sec:leakage}).  This
section sets out the comparison that would settle the question.  The
protocol is written so that the GBM \emph{can fail}: if a challenger
wins under the pre-agreed decision rule, the challenger is adopted.

\subsection{Why Random Forest Is the Natural Challenger}

Both models are tree ensembles on the same four features, so they
differ only in how the ensemble is built:

\begin{center}
\begin{tabular}{p{3.2cm}p{5.2cm}p{5.2cm}}
\toprule
& \textbf{GBM (champion)} & \textbf{Random Forest (challenger)} \\
\midrule
Ensemble     & Sequential, each tree fits the residual of the last
             & Independent deep trees on bootstrap samples, averaged \\
Error reduced & Mainly bias & Mainly variance \\
Tuning       & Sensitive (learning rate, depth, trees)
             & Robust to defaults \\
Probabilities & Optimises log loss directly; usually well calibrated
             & Averaged votes; pulled away from 0 and 1, often needs
               recalibration \\
Seed sensitivity & Low--moderate (\texttt{subsample=0.8})
             & Low with many trees \\
\bottomrule
\end{tabular}
\end{center}

\noindent Because the classifier output is used directly as a price
($C_h = N \times p_h \times d$), calibration matters more than
ranking.  This is the main reason to expect the GBM to win, but it has
to be shown, not assumed.

\subsection{Candidate Set}

\begin{enumerate}
  \item \textbf{GBM} --- production configuration
        (\texttt{n\_estimators=300}, \texttt{max\_depth=3},
        \texttt{learning\_rate=0.05}, \texttt{subsample=0.8},
        \texttt{min\_samples\_leaf=50}).
  \item \textbf{GBM, tuned} --- the same model given the same tuning
        budget as the challengers.  If it beats the production GBM,
        the production settings, not the model family, need changing.
  \item \textbf{Random Forest} --- \texttt{n\_estimators=500},
        \texttt{max\_features=2}, \texttt{min\_samples\_leaf} tuned.
  \item \textbf{Random Forest, calibrated} --- the same, wrapped in
        isotonic or Platt calibration fitted on an inner fold.  This
        stops RF losing only because of calibration that a
        one-line fix would remove.
  \item \textbf{Logistic regression} on the four features --- a
        transparent baseline.
  \item \textbf{Flood polynomial} --- the production fallback, as the
        floor every model must beat.
\end{enumerate}

\noindent Every tree model gets the same tuning budget (for example
30 random-search trials) inside the training folds, so the GBM does
not win only because it has been tuned more.

\subsection{Prerequisite: Remove Storm Leakage}
\label{sec:leakage}

Each storm contributes 168 hourly rows, and the label
(\emph{future maximum level} $\geq w_{\text{severe}}$) is shared by
adjacent hours.  The current trainer uses a row-level stratified
\texttt{train\_test\_split}, so hours from the same storm appear in
both training and test sets.  The reported AUC therefore measures
interpolation within storms already seen, not prediction on new
storms.  This leakage favours deep, low-bias learners and so would
flatter the Random Forest more than the GBM.

\noindent Before any comparison:
\begin{itemize}
  \item Carry a \texttt{storm\_id} with every training vector.
  \item Split by storm using \texttt{GroupKFold} (5 folds), so every
        storm is entirely in train or entirely in test.
  \item Add near-miss synthetic vectors to the \emph{training} folds
        only; score on real storm vectors only.
\end{itemize}

\subsection{Accuracy Metrics}

Scored on held-out storms, per gauge and pooled:
\begin{description}[leftmargin=2em]
  \item[Primary] Brier score and log loss --- proper scoring rules
        that reward correct probabilities, not just correct ranking.
  \item[Calibration] Reliability diagram, expected calibration error,
        calibration slope and intercept.
  \item[Discrimination] AUC (secondary, kept for continuity with
        the existing thresholds).
  \item[Business] Stress P\&L error on held-out storms: the hourly
        $\sum_h N p_h d$ path compared against the realised
        knock-out outcome of that storm.
\end{description}

\subsection{Stability Tests}

``Most stable'' is tested in five ways.  Each produces a number per
model and per gauge.

\begin{enumerate}
  \item \textbf{Seed stability.}  Retrain each model with 20 random
        seeds on the same data.  Measure the standard deviation of
        $p$ on a fixed grid (the probability surface of
        Section~4) and of peak stress P\&L for a reference storm set.
  \item \textbf{Data stability.}  Retrain on 50 bootstrap resamples
        \emph{of storms} (not of rows).  Report the spread of $p$ on
        the fixed grid and the spread of the held-out Brier score.
  \item \textbf{Catalogue stability.}  Train on one generation of the
        storm catalogue and score on a regenerated catalogue
        (a new seed or new storm-control parameters).  This is the
        retraining trigger in practice, so it is the most
        decision-relevant stability test.  The benchmark tool does not
        regenerate the catalogue itself: run it once on each catalogue
        generation and compare the two result files.
  \item \textbf{Shape stability.}  $P(\text{flood})$ should not fall
        as the water level rises.  Count monotonicity violations in
        \texttt{log\_h\_s} across the grid, and measure how rough the
        finite-difference sensitivity $\partial p / \partial w$ is.
        Step-like surfaces give noisy stress Greeks.
  \item \textbf{Cross-gauge consistency.}  The distribution of
        per-gauge metrics: a model that is excellent on most gauges
        but fails badly on a few is less stable than one that is
        uniformly good.
\end{enumerate}

\subsection{Statistical Comparison}

\begin{itemize}
  \item Per gauge: paired difference in Brier score (GBM minus
        challenger) on the same held-out storms, with a 95\%
        confidence interval from a storm-level bootstrap.
  \item Across gauges: Wilcoxon signed-rank test on the per-gauge
        Brier differences.
  \item Stability metrics: report ratios (challenger dispersion /
        GBM dispersion) with bootstrap intervals.
\end{itemize}

\subsection{Decision Rule}

Agreed before the results are seen.  The GBM is retained if all of
the following hold:
\begin{enumerate}
  \item Its pooled Brier score is no worse than the best challenger's
        by more than a set tolerance (for example 2\%), and the
        Wilcoxon test does not show the challenger significantly
        better.
  \item Its seed, bootstrap and catalogue dispersion are no greater
        than the challenger's (within 0.002 in probability s.d.).
  \item It has no more monotonicity violations than the challenger
        (within 0.5 percentage points of the violation rate).
\end{enumerate}

\noindent If a challenger is better on accuracy but worse on
stability (or the reverse), the trade-off is recorded and the choice
is escalated for model-risk sign-off rather than decided by the
model owner.  If the best challenger wins on all three, it replaces
the GBM.  With fewer than six gauges the Wilcoxon test has no power, so
``significantly better'' then means every gauge's storm-bootstrap
interval favours the challenger.

\subsection{Running the Benchmark}

The comparison is implemented in the benchmark tool:
\begin{description}[leftmargin=2em]
  \item[Code] \texttt{tools/flood\_classifier\_challenger/}
  \item[Settings] \texttt{config/models/\_flood\_challenger.py}
\end{description}
\noindent It builds each gauge's training set with the production
vector builder, so it scores exactly the features and labels
production trains on:

\begin{verbatim}
python -m tools.flood_classifier_challenger \
    --gauges all
python -m tools.flood_classifier_challenger \
    --gauges G1,G2 --seeds 20 --bootstraps 50
python -m tools.flood_classifier_challenger \
    --synthetic 3     # no storm data needed
\end{verbatim}

\noindent Defaults are 1{,}000 storms per gauge, 5 storm folds,
6 tuning trials, 10 seeds and 20 storm bootstraps; a sign-off run
uses 20 seeds and 50 bootstraps as above.  Fits run in parallel on
every core.  The run writes \texttt{challenger/challenger\_results.json}
(every metric, the tuned parameters per fold and the probability
surfaces) and the table below.

\subsection{Results}

\IfFileExists{challenger/challenger_results.tex}{%
  % Generated by tools/flood_classifier_challenger -- do not edit by hand.
\noindent\textbf{Run:} 2026-10-30, source \texttt{catalogue thames from PhysicalRisk-testdata}, 6 gauges scored (0 skipped), 500 storms per gauge, 5 storm folds, 6 tuning trials, 10 seeds, 20 storm bootstraps.

\begin{center}\small
\begin{tabular}{lccccccc}
\toprule
\textbf{Model} & \textbf{Brier} & \textbf{Log loss} & \textbf{AUC} & \textbf{ECE} & \textbf{Seed s.d.} & \textbf{Boot.\ s.d.} & \textbf{Mono.\ viol.} \\
\midrule
GBM (production) & 0.0204 & 0.0688 & 0.989 & 0.0078 & 0.0090 & 0.0353 & 8.3\% \\
GBM (tuned) & 0.0191 & 0.0644 & 0.990 & 0.0051 & 0.0000 & 0.0495 & 8.9\% \\
Random Forest & 0.0185 & 0.0631 & 0.990 & 0.0054 & 0.0052 & 0.0354 & 11.2\% \\
Random Forest (calib.) & 0.0199 & 0.0752 & 0.988 & 0.0061 & 0.0420 & 0.0610 & 5.3\% \\
Logistic regression & 0.0378 & 0.1438 & 0.921 & 0.0244 & 0.0000 & 0.0175 & 0.0\% \\
Flood polynomial & 0.0203 & 0.0822 & 0.979 & 0.0074 & -- & -- & 11.5\% \\
\bottomrule
\end{tabular}
\end{center}

\noindent\textbf{Decision:} best challenger by pooled Brier is Random Forest (GBM Brier 0.0204 vs 0.0185, gap 10.03\%; Wilcoxon $p = 0.016$).  Accuracy test failed, stability test failed, shape test passed: the result is \textbf{mixed} and is escalated for model-risk sign-off.

\begin{center}\small
\begin{longtable}{lrrccc}
\toprule
\textbf{Gauge} & \textbf{Storms} & \textbf{Flood rate} & \textbf{GBM Brier} & \textbf{Best chall.} & \textbf{$\Delta$ (95\% CI)} \\
\midrule\endhead
GAUGE-ef7f75ec & 500 & 12.6\% & 0.0297 & 0.0270 & +0.0027 [+0.0015, +0.0039] \\
GAUGE-066fc60e & 500 & 7.3\% & 0.0188 & 0.0166 & +0.0022 [+0.0010, +0.0034] \\
GAUGE-bc801f07 & 500 & 3.6\% & 0.0111 & 0.0102 & +0.0009 [+0.0001, +0.0018] \\
GAUGE-3324d948 & 500 & 7.3\% & 0.0210 & 0.0189 & +0.0022 [+0.0010, +0.0032] \\
GAUGE-acc1b9cc & 500 & 6.5\% & 0.0208 & 0.0197 & +0.0011 [+0.0001, +0.0022] \\
GAUGE-a03e456e & 500 & 7.1\% & 0.0208 & 0.0187 & +0.0021 [+0.0008, +0.0033] \\
\bottomrule
\end{longtable}
\end{center}
\noindent $\Delta$ is GBM minus challenger per-storm Brier; positive favours the challenger.
%

\medskip\noindent\textbf{Reading these results.}  This run is
illustrative, not sign-off evidence: it uses the first six gauges of
the 18 in the test portfolio, at 500 storms each and the default
seed and bootstrap counts.  The tuned GBM's zero seed s.d.\ does not
show better stability: tuning chose \texttt{subsample} $= 1$ for the
fold used in the stability test on all six gauges, so the fit has no
randomness for the seed to act on.  Its bootstrap s.d., which
measures sensitivity to the storms sampled, is the higher of the
GBM variants and above the Random Forest's.  Tuned parameters also
differ between folds, so its accuracy is not that of one fixed
setting.%
}{%
\noindent The comparison has not yet been run on the storm catalogue.
Until it is, the choice of GBM rests on modelling judgement and
in-sample AUC, and should be recorded as such.}


% =========================================================================
\section{Parameter Reference}
\label{ch:reference}
% =========================================================================

\begin{longtable}{p{4.8cm}rp{3.2cm}p{3.2cm}}
\toprule
\textbf{Parameter} & \textbf{Default} & \textbf{Increase} & \textbf{Decrease} \\
\midrule
\endfirsthead
\toprule
\textbf{Parameter} & \textbf{Default} & \textbf{Increase} & \textbf{Decrease} \\
\midrule
\endhead
\texttt{STRESS\_STORMS\_MIN\_COUNT} & 50 &
  Excludes gauges with few storms &
  Includes more gauges \\[4pt]
\texttt{STORM\_DEFAULT\_DURATION} & 168\,h &
  Longer fallback hydrograph &
  Shorter fallback \\[4pt]
\texttt{DEFAULT\_PEAK\_POSITION} & 0.5 &
  Late peak (delayed climax) &
  Early peak (front-loaded) \\[4pt]
\texttt{LOG\_END} & 168.0 &
  Classifier feature scaling &
  Narrower time window \\[4pt]
\texttt{LOG\_EPS} & $10^{-8}$ &
  Less numerical sensitivity &
  More sensitivity at zero \\[4pt]
\texttt{BUMP\_1BP} & 0.0001 &
  Larger sensitivity step &
  Finer stress resolution \\
\bottomrule
\end{longtable}


% =========================================================================
\section*{Document History}
% =========================================================================

\begin{tabular}{llll}
\toprule
\textbf{Version} & \textbf{Date} & \textbf{Author} & \textbf{Description} \\
\midrule
1.0 & 10-Apr-2026 & MKM Research Labs & Initial release \\
1.1 & 29-Oct-2026 & MKM Research Labs & Added GBM challenger comparison \\
\bottomrule
\end{tabular}

\end{document}
